\chapter{Fields}

A \emph{field} \((\field, \oplus, \odot)\) is a commutative ring where every nonzero element has a \emph{multiplicative inverse}:

\[
	(\field \setminus \{0\}, \odot) \text{ is an Abelian group}.
\]

Fields like \(\Q\), \(\R\), and \(\C\) allow division (except by zero), enabling the full range of arithmetic operations.

\section{Field Properties}

\begin{enumerate}

	\item \emph{Neutral Elements are unique}: There exists exactly one additive identity \(0\) and one multiplicative identity \(1\) in a field.

	\item \emph{Inverses are unique}: For each \(a \in \field\), there exists exactly one additive inverse \(-a\) and, for \(a \ne 0\), exactly one multiplicative inverse \(a^{-1}\).

	\item \emph{Any element multiplied by zero is zero}: For any \(a \in \field\), \(a \cdot 0 = 0\).

	\item \emph{Factorization of \(-1\)}: For any \(a \in \field\), \(-a= (-1)a\).

	\item \emph{Distribution of negative signs for multiplication}: For any \(a, b \in \field\), \(-(ab) = (-a)b = a(-b)\).

	\item \emph{Minus times minus is plus}: For any \(a, b \in \field\), \((-a)(-b) = ab\).

	\item \emph{Uniqueness of differences}: For any \(a, b \in \field\), there exists exactly one \(x \in \field\) such that \(a + x = b\), denoted as \(x = b - a\).

	\item \emph{Distribution of negative signs}: For any \(a, b \in \field\), \(-(a + b) = (-a) + (-b)\).

	\item \emph{Zero Product}: For any \(a, b \in \field\), if \(ab = 0\), then either \(a = 0\) or \(b = 0\).

\end{enumerate}

\textbf{Proof I:}

For \(a + 0 = a\). Now let us say \(0' \in \field\), such that \(a + 0' = a\). Then we have for \(a = 0\):

\[
	0 + 0' = 0 = 0' + 0 \implies 0' = 0
\]

For multiplicative neutral element, it is the same proof, but with the multiplicative operation. Hence
both neutral elements are unique.

\QED

\textbf{Proof II:}

Let \(a + -a = 0\). Now let us say \(a' \in \field\) such that \(a + a' = 0\). Then we have

\begin{align*}
	a + -a        & = a + a'        \\
	(a + -a) + -a & = a + (a' + -a) \\
	0 + -a        & = (a + -a) + a' \\
	0 + -a        & = 0 + a'        \\
	-a            & = a'
\end{align*}

Hence the additive inverse is unique. The proof for the multiplicative inverse is analogous.

\QED

\textbf{Proof III:}

Let \(a \in \field\) and \(0 \in \field\) be the additive identity. Then we have

\begin{align*}
	a \cdot 0 & = (a \cdot 0) + 0                             \\
	          & = (a \cdot 0) + ((a \cdot 0) + - (a \cdot 0)) \\
	          & = ((a \cdot 0) + (a \cdot 0)) + - (a \cdot 0) \\
	          & = a \cdot(0 + 0) + - (a \cdot 0)              \\
	          & = a \cdot 0 + - (a \cdot 0)                   \\
	          & = 0
\end{align*}

\QED

\textbf{Proof IV:}

Let \(-a \in \field\) then,

\begin{align*}
	(1 + -1)       & = 0  \\
	x + (-1)x      & = 0  \\
	x + (-1)x + -x & = -x \\
	x + -x + (-1)x & = -x \\
	0 + (-1)x      & = -x \\
	(-1)x          & = -x
\end{align*}

\QED

\textbf{Proof V:}

Let \(a, b \in \field\) then,

\[
	-ab = ((-1)a)b \iff (-1)(ab) = -(ab)
\]

\QED

\textbf{Proof VI:}

Let \(a, b \in \field\) then, we know that \(ab + -(ab) = 0\). Therefore,

\begin{align*}
	ab + -(ab)            & = 0                                                \\
	ab + -(ab) + -(-(ab)) & = 0 + -(-(ab))                                     \\
	ab                    & = -(-(ab)) = (-1)(-(ab)) = (-1)(-1)(ab) = (-a)(-b)
\end{align*}

\QED

\textbf{Proof VII:}

There exists exactly one \(x \in \field\) such that \(a + x = b\), \(x = b + (-a)\) also written
\(x = b - a\).

Let \(x, x' \in \field\) such that the property is satified then,

\[
	x = b - a \land x' = b - a \implies x = x'.
\]

\QED

A special case of this property is when \(b = a\), which gives us that

\[
	a + x = a \land a + x' = a \iff a + 0 = a \land a - 0 = a \implies 0 = -0.
\]

\textbf{Proof VIII:}

For any \(a, b \in \field\) the following holds \(-(a + b)  = (-a) + (-b)\).

Note that \(-(a + b)\) satisfies

\[
	(a + b) - (a + b) = 0
\]

Now note that

\begin{align*}
	(a + b) - a - b & =  (b + a) - a - b \\
	                & = b + (a - a) - b  \\
	                & = b - b = 0
\end{align*}

Hence, \(-a - b\) also satisfies the same difference. By our previous property we know that
difference are unique hence, \(-(a + b) = (-a) + (-b)\).

\QED


